% Calorimètre (vase de Dewar) : méthode des mélanges (L'essentiel, chapitre 19)
\begin{tikzpicture}[x=1cm,y=1cm]
% double paroi
\draw[otInk!70,line width=0.8pt,rounded corners=2mm] (0,0) rectangle (2.6,2.8);
\draw[otInk!70,line width=0.6pt,rounded corners=1.5mm] (0.18,0.18) rectangle (2.42,2.6);
\fill[otMuted!15] (0.02,0.02) rectangle (0.16,2.6) (2.44,0.02) rectangle (2.58,2.6) (0.16,0.02) rectangle (2.44,0.16);
\fill[pcEau] (0.2,0.2) rectangle (2.4,1.9);
\draw[otNavy!45,line width=0.5pt] (0.2,1.9) -- (2.4,1.9);
% couvercle
\fill[otMuted!55] (-0.1,2.8) rectangle (2.7,3.05);
% thermomètre
\fill[white] (0.75,0.5) rectangle (0.9,3.7);
\draw[otInk!70,line width=0.5pt] (0.75,0.5) rectangle (0.9,3.7);
\fill[otRed!80!black] (0.825,0.5) circle (0.12);
\fill[otRed!80!black] (0.79,0.5) rectangle (0.86,2.3);
% agitateur
\draw[otInk!70,line width=0.9pt] (1.75,0.45) -- (1.75,3.6);
\draw[otInk!70,line width=0.9pt] (1.45,0.45) -- (2.05,0.45);
\draw[otInk!70,line width=0.9pt] (1.55,3.6) -- (1.95,3.6);
% légendes
\draw[otMuted,line width=0.3pt] (0.9,3.3) -- (3.3,3.3) node[pcLegende,anchor=west] {thermomètre};
\draw[otMuted,line width=0.3pt] (1.95,3.6) -- (3.3,3.65) node[pcLegende,anchor=west] {agitateur};
\draw[otMuted,line width=0.3pt] (2.51,1.4) -- (3.3,1.4) node[pcLegende,anchor=west,align=left] {double paroi isolante\\ (vide entre les parois)};
\draw[otMuted,line width=0.3pt] (2.0,1.0) -- (3.3,0.6) node[pcLegende,anchor=west] {eau et corps étudié};
\end{tikzpicture}
\par\smallskip\legende{Le calorimètre limite les transferts thermiques avec l'extérieur : les corps qu'il contient n'échangent de l'énergie qu'entre eux.}
